\section{Conclusion}


Jin demonstrates that joint embedding architectures can unify multiple modalities into a single, coherent representation space. The combination of modality-specific encoders, diffusion transformer backbone with MoE, and hierarchical z-JEPA enables state-of-the-art performance on multimodal understanding tasks while maintaining computational efficiency. The architecture's extensibility opens new possibilities for AI systems that perceive and interact with the world across all sensory modalities.

\bibliographystyle{plain}
\begin{thebibliography}{10}

\bibitem{jepa2022}
Y. LeCun. A Path Towards Autonomous Machine Intelligence. 2022.

\bibitem{moe2017}
N. Shazeer et al. Outrageously Large Neural Networks: The Sparsely-Gated Mixture-of-Experts Layer. ICLR, 2017.

\bibitem{ddpm2020}
J. Ho et al. Denoising Diffusion Probabilistic Models. NeurIPS, 2020.

\bibitem{dit2023}
W. Peebles, S. Xie. Scalable Diffusion Models with Transformers. ICCV, 2023.

\bibitem{vit2020}
A. Dosovitskiy et al. An Image is Worth 16x16 Words. ICLR, 2021.

\bibitem{ijepa2023}
M. Assran et al. Self-Supervised Learning from Images with a Joint-Embedding Predictive Architecture. CVPR, 2023.

\bibitem{vjepa2024}
A. Bardes et al. V-JEPA: Latent Video Prediction for Visual Representation Learning. 2024.

\bibitem{clip2021}
A. Radford et al. Learning Transferable Visual Models From Natural Language Supervision. ICML, 2021.

\end{thebibliography}

